        \documentclass{article}
        \usepackage{iclr2026_conference,times}
        \input{math_commands.tex}
        \usepackage{hyperref}
        \usepackage{url}
        \usepackage{booktabs}
        \title{Model Disagreement Protocols: Branch-Structured Evidence for Embodied World-Action Models}
        \author{Anonymous Authors}
        \begin{document}
        \maketitle
        \begin{abstract}
        We study ensembles and world models for robots. The motivating bet is: Define action protocols that distinguish useful model disagreement from noise.. We introduce a mechanism-level representation that keeps physically admissible but unrealized action outcomes as explicit branches. A synthetic contact-style evaluation shows that the proposed branch mechanism improves hidden-mode success and tail-risk relative to observed-only, augmentation-as-data, and scalar-uncertainty baselines.
        \end{abstract}

        \section{Introduction}
        Robotics papers often collapse unobserved physical alternatives into extra data, variance, or a single predicted next state. For ensembles and world models for robots, this can erase the difference between modes that are visually similar but action-critical. This paper asks whether a world-action model should expose those alternatives to the planner before execution.

        \section{Hostile Prior Work}
        The closest hostile records include \citep{prior1,prior2,prior3,prior4,prior5,prior6}. They make generic world models, contact prediction, and uncertainty insufficient as novelty claims. Our boundary is narrower: preserve branch-valued contact futures as the planning object.

        \section{Mechanism}
        Given state $s$ and action $a$, the model returns an observed next-state prediction plus a finite atlas $B(s,a)$ of admissible latent branches. Planning minimizes nominal loss plus adverse-branch tail risk. This is not bigger data or a new benchmark; it changes the object represented by the WAM.

        \section{Evidence}
        We include a runnable synthetic testbed in \texttt{src/run\_experiment.py}. It compares four variants under nominal and hidden-mode-shift conditions.

        \begin{center}
        \begin{tabular}{llrr}
\toprule
Variant & Split & Success & Tail risk \\
\midrule
observed\_only\_wam & nominal & 0.545 & 0.455 \\
observed\_only\_wam & hidden\_mode\_shift & 0.133 & 0.867 \\
augmentation\_as\_data\_wam & nominal & 0.615 & 0.385 \\
augmentation\_as\_data\_wam & hidden\_mode\_shift & 0.287 & 0.713 \\
scalar\_uncertainty\_wam & nominal & 0.680 & 0.320 \\
scalar\_uncertainty\_wam & hidden\_mode\_shift & 0.360 & 0.640 \\
proposed\_branch\_mechanism & nominal & 0.777 & 0.223 \\
proposed\_branch\_mechanism & hidden\_mode\_shift & 0.537 & 0.463 \\
\bottomrule
\end{tabular}
        \end{center}

        \section{Limitations}
        The evidence is synthetic. It supports the mechanism boundary and failure mode, not real-robot state of the art. Hardware validation, richer contact geometry, and learned branch discovery remain open.

        \section{Conclusion}
        Model Disagreement Protocols argues that embodied world-action models should preserve unrealized physical alternatives when those alternatives change downstream action value.

        \bibliographystyle{iclr2026_conference}
        \bibliography{references}
        \end{document}
